\documentclass{article}

\usepackage{../paper/neurips_2026}

\usepackage[utf8]{inputenc}
\usepackage[T1]{fontenc}
\usepackage{hyperref}
\usepackage{url}
\usepackage{booktabs}
\usepackage{amsfonts}
\usepackage{nicefrac}
\usepackage{microtype}
\usepackage{xcolor}
\usepackage{amsmath,amssymb,amsthm,mathtools}
\usepackage{bm}
\usepackage{enumitem}
\usepackage[nameinlink,noabbrev]{cleveref}

\newtheorem{theorem}{Theorem}
\newtheorem{proposition}[theorem]{Proposition}
\newtheorem{corollary}[theorem]{Corollary}
\newtheorem{lemma}[theorem]{Lemma}
\theoremstyle{definition}
\newtheorem{definition}[theorem]{Definition}
\newtheorem{remark}[theorem]{Remark}

\newcommand{\R}{\mathbb{R}}
\newcommand{\E}{\mathbb{E}}
\newcommand{\PP}{\mathcal{P}}
\newcommand{\W}{W_1}
\newcommand{\dd}{\,\mathrm{d}}

\title{A Three-Layer BatchEnsemble-Like Proxy for the Large-$K$ Regime}

\author{Anonymous Authors}

\begin{document}

\maketitle

\begin{abstract}
This note records a semifaithful three-layer BatchEnsemble-like proxy for the large-$K$ regime and explains why it is both more faithful to practical parameter-efficient ensembling and substantially harder than the two-layer case. The key structural point survives: at fixed width, parameter sharing still leaves the pair consisting of the shared state and the empirical head law as an exact closed state variable. The new ingredient is a shared hidden-to-hidden matrix of size $m\times m$, which is precisely where BatchEnsemble-style parameter efficiency becomes substantial. For the concrete three-layer proxy we derive explicit gradient and aggregate-norm bounds. These bounds reveal the main obstruction: unlike in two layers, the natural majorant system is already Riccati-type at the upper hidden layer and becomes higher-order at the first layer. Thus the same bounded-support localization argument used in the rigorous two-layer theorem no longer closes automatically.
\end{abstract}

\section{Motivation}

The two-layer proxy from the companion note is a useful first rigorous target, but it does not yet contain a shared hidden-to-hidden matrix of size $m\times m$. In a practical efficient ensemble, this is exactly where parameter sharing is most valuable: replacing $K$ independent $m\times m$ matrices by one shared matrix together with $O(Km)$ head-specific modulation vectors is the main source of parameter savings. A three-layer reduced proxy is therefore the first setting in which the parameter-efficiency mechanism of BatchEnsemble becomes genuinely visible.

\section{A Semifaithful Three-Layer Proxy}

Fix an input dimension $d$ and a hidden width $m$. We consider
\begin{align}
  h_{1,\alpha}(x)
  &:=
  s_{1,\alpha}\odot
  \sigma\bigl(W_1(r_{1,\alpha}\odot x)+b_1\bigr)
  \in\R^m,
  \label{eq:h1}
  \\
  h_{2,\alpha}(x)
  &:=
  s_{2,\alpha}\odot
  \sigma\bigl(W_2(r_{2,\alpha}\odot h_{1,\alpha}(x))+b_2\bigr)
  \in\R^m,
  \label{eq:h2}
  \\
  f_\alpha(x)
  &:=
  \frac1m a^\top h_{2,\alpha}(x),
  \qquad
  \alpha=1,\dots,K.
  \label{eq:f3}
\end{align}
Here
\[
  a\in\R^m,\quad
  W_1\in\R^{m\times d},\quad
  b_1\in\R^m,\quad
  W_2\in\R^{m\times m},\quad
  b_2\in\R^m
\]
are shared, while each head carries
\[
  s_{1,\alpha}\in\R^m,\quad
  r_{1,\alpha}\in\R^d,\quad
  s_{2,\alpha}\in\R^m,\quad
  r_{2,\alpha}\in\R^m.
\]
Thus the shared parameter count is $O(md+m^2)$, while the per-head parameter count is only $O(d+3m)$.

We collect the shared and head states as
\[
  q:=(a,W_1,b_1,W_2,b_2)\in\mathcal Q
  :=
  \R^m\times\R^{m\times d}\times\R^m\times\R^{m\times m}\times\R^m,
\]
\[
  z:=(s_1,r_1,s_2,r_2)\in\mathcal Z
  :=
  \R^m\times\R^d\times\R^m\times\R^m.
\]
Writing $f(x;q,z)$ for the output \eqref{eq:f3} of a generic head state against a shared state, the head-wise population risk is
\[
  \mathcal R_K
  =
  \E_{(x,y)}
  \Bigl[
    \frac1K\sum_{\alpha=1}^K \ell(f(x;q,z_\alpha),y)
  \Bigr].
\]
As in the two-layer note, the induced shared and head drifts are
\[
  \mathcal F(q,z)
  :=
  -\E
  \bigl[
    \partial_1\ell(f(x;q,z),y)\,\nabla_q f(x;q,z)
  \bigr],
\]
\[
  \mathcal G(q,z)
  :=
  -\E
  \bigl[
    \partial_1\ell(f(x;q,z),y)\,\nabla_z f(x;q,z)
  \bigr].
\]

\begin{proposition}[Exact empirical closure]
\label{prop:3layer-closure}
For every fixed $K$ and every head-learning scale $\lambda_K\ge 0$, the finite-$K$ system
\[
  \dot q_t^K
  =
  \frac1K\sum_{\alpha=1}^K \mathcal F(q_t^K,z_{\alpha,t}^K),
  \qquad
  \dot z_{\alpha,t}^K
  =
  \lambda_K\mathcal G(q_t^K,z_{\alpha,t}^K)
\]
is exactly closed at the level of $(q_t^K,\nu_t^K)$ with
\[
  \nu_t^K:=\frac1K\sum_{\alpha=1}^K \delta_{z_{\alpha,t}^K}.
\]
In particular,
\[
  \dot q_t^K
  =
  \int_{\mathcal Z}\mathcal F(q_t^K,z)\,\nu_t^K(\dd z),
\]
and
\[
  \partial_t\nu_t^K
  +
  \lambda_K\nabla_z\cdot\bigl(\nu_t^K \mathcal G(q_t^K,z)\bigr)
  =
  0
\]
in the distributional sense.
\end{proposition}

\begin{proof}
The proof is identical to the two-layer case. It uses only the factorization
\[
  f_\alpha(x)=f(x;q,z_\alpha)
\]
and permutation invariance in the head labels.
\end{proof}

\section{Concrete Gradient Bounds}

Assume throughout this section that
\[
  \|x\|\le B_X
  \qquad
  \mathsf P\text{-almost surely},
\]
\[
  \|\sigma\|_\infty\le M_0,
  \qquad
  \|\sigma'\|_\infty\le M_1,
  \qquad
  |\partial_1\ell(u,y)|\le L_\ell.
\]

\begin{proposition}[Three-layer gradient bounds]
\label{prop:3layer-grads}
For the model \eqref{eq:h1}--\eqref{eq:f3}, the following bounds hold for every $(x,q,z)$:
\begin{align}
  \|\nabla_a f\|
  &\le
  \frac{M_0}{m}\|s_2\|,
  \label{eq:3g-a}
  \\
  \|\nabla_{s_2} f\|
  &\le
  \frac{M_0}{m}\|a\|,
  \label{eq:3g-s2}
  \\
  \|\nabla_{b_2} f\|
  &\le
  \frac{M_1}{m}\|a\|\,\|s_2\|,
  \label{eq:3g-b2}
  \\
  \|\nabla_{W_2} f\|_F
  &\le
  \frac{M_0M_1}{m}
  \|a\|\,\|s_2\|\,\|r_2\|\,\|s_1\|,
  \label{eq:3g-W2}
  \\
  \|\nabla_{r_2} f\|
  &\le
  \frac{M_0M_1}{m}
  \|a\|\,\|s_2\|\,\|W_2\|_F\,\|s_1\|,
  \label{eq:3g-r2}
  \\
  \|\nabla_{s_1} f\|
  &\le
  \frac{M_0M_1}{m}
  \|a\|\,\|s_2\|\,\|r_2\|\,\|W_2\|_F,
  \label{eq:3g-s1}
  \\
  \|\nabla_{b_1} f\|
  &\le
  \frac{M_0M_1^2}{m}
  \|a\|\,\|s_2\|\,\|r_2\|\,\|W_2\|_F\,\|s_1\|,
  \label{eq:3g-b1}
  \\
  \|\nabla_{W_1} f\|_F
  &\le
  \frac{M_0M_1^2B_X}{m}
  \|a\|\,\|s_2\|\,\|r_2\|\,\|W_2\|_F\,\|s_1\|\,\|r_1\|,
  \label{eq:3g-W1}
  \\
  \|\nabla_{r_1} f\|
  &\le
  \frac{M_0M_1^2B_X}{m}
  \|a\|\,\|s_2\|\,\|r_2\|\,\|W_2\|_F\,\|s_1\|\,\|W_1\|_F.
  \label{eq:3g-r1}
\end{align}
\end{proposition}

\begin{proof}
The first three bounds are immediate from the explicit formulas
\[
  \nabla_a f
  =
  \frac1m h_2,
  \qquad
  \nabla_{s_2} f
  =
  \frac1m a\odot \sigma(u_2),
  \qquad
  \nabla_{b_2} f
  =
  \frac1m a\odot s_2\odot \sigma'(u_2).
\]
For \eqref{eq:3g-W2}, the $j$th row of $\nabla_{W_2}f$ is
\[
  \frac{a_j s_{2,j}}{m}
  \sigma'(u_{2,j})
  (r_2\odot h_1)^\top,
\]
so
\[
  \|\nabla_{W_2}f\|_F
  \le
  \frac{M_1}{m}\|a\odot s_2\|\,\|r_2\odot h_1\|
  \le
  \frac{M_0M_1}{m}
  \|a\|\,\|s_2\|\,\|r_2\|\,\|s_1\|.
\]
For \eqref{eq:3g-r2}, write
\[
  v_2
  :=
  \frac1m W_2^\top(a\odot s_2\odot \sigma'(u_2)).
\]
Then
\[
  \nabla_{r_2}f
  =
  h_1\odot v_2,
\]
so
\[
  \|\nabla_{r_2}f\|
  \le
  \|h_1\|_\infty \|v_2\|
  \le
  M_0\|s_1\|
  \frac{M_1}{m}\|W_2\|_F \|a\odot s_2\|,
\]
which gives \eqref{eq:3g-r2}. Likewise,
\[
  \nabla_{s_1}f
  =
  \sigma(u_1)\odot (r_2\odot v_2),
\]
which yields \eqref{eq:3g-s1}. Finally,
\[
  \nabla_{b_1}f
  =
  s_1\odot \sigma'(u_1)\odot (r_2\odot v_2),
\]
and the bounds for $\nabla_{W_1}f$ and $\nabla_{r_1}f$ follow by one more chain-rule step, using $\|r_1\odot x\|\le B_X\|r_1\|$ and $\|W_1^\top u\|\le \|W_1\|_F\|u\|$.
\end{proof}

\section{What Breaks Relative to Two Layers}

For the finite-$K$ dynamics, define the shared and empirical head norms
\[
  A_t:=\|a_t\|,
  \qquad
  W_{1,t}:=\|W_{1,t}\|_F,
  \qquad
  W_{2,t}:=\|W_{2,t}\|_F,
\]
\[
  S_{1,t}
  :=
  \Bigl(
    \frac1K\sum_{\alpha=1}^K \|s_{1,\alpha,t}\|^2
  \Bigr)^{1/2},
  \quad
  R_{1,t}
  :=
  \Bigl(
    \frac1K\sum_{\alpha=1}^K \|r_{1,\alpha,t}\|^2
  \Bigr)^{1/2},
\]
\[
  S_{2,t}
  :=
  \Bigl(
    \frac1K\sum_{\alpha=1}^K \|s_{2,\alpha,t}\|^2
  \Bigr)^{1/2},
  \quad
  R_{2,t}
  :=
  \Bigl(
    \frac1K\sum_{\alpha=1}^K \|r_{2,\alpha,t}\|^2
  \Bigr)^{1/2}.
\]
Applying \cref{prop:3layer-grads} exactly as in the two-layer proof gives a closed system of differential inequalities of the form
\begin{align}
  \dot A_t
  &\lesssim
  S_{2,t},
  \label{eq:3d-A}
  \\
  \dot S_{2,t}
  &\lesssim
  A_t,
  \label{eq:3d-S2}
  \\
  \dot W_{2,t}
  &\lesssim
  A_t S_{2,t} S_{1,t} R_{2,t},
  \label{eq:3d-W2}
  \\
  \dot R_{2,t}
  &\lesssim
  A_t S_{2,t} S_{1,t} W_{2,t},
  \label{eq:3d-R2}
  \\
  \dot S_{1,t}
  &\lesssim
  A_t S_{2,t} R_{2,t} W_{2,t},
  \label{eq:3d-S1}
  \\
  \dot W_{1,t}
  &\lesssim
  A_t S_{2,t} R_{2,t} W_{2,t} S_{1,t} R_{1,t},
  \label{eq:3d-W1}
  \\
  \dot R_{1,t}
  &\lesssim
  A_t S_{2,t} R_{2,t} W_{2,t} S_{1,t} W_{1,t}.
  \label{eq:3d-R1}
\end{align}

\begin{remark}[Upper and lower layers behave differently]
\label{rem:3layer-hierarchy}
The top layer still closes benignly:
\[
  V_t:=A_t+S_{2,t}
  \quad\Longrightarrow\quad
  \dot V_t \lesssim V_t,
\]
so $V_t$ grows at most exponentially. The first nontrivial obstruction appears at the hidden-to-hidden layer. If
\[
  U_t:=1+S_{1,t}+R_{2,t}+W_{2,t},
\]
then \eqref{eq:3d-W2}--\eqref{eq:3d-S1} imply only
\[
  \dot U_t \lesssim C_T U_t^2,
\]
because the bounded top-layer variable $V_t$ enters only through a finite-horizon constant $C_T$. This is already Riccati-type. The first layer is worse still: with
\[
  L_t:=1+W_{1,t}+R_{1,t},
\]
\eqref{eq:3d-W1}--\eqref{eq:3d-R1} give only
\[
  \dot L_t \lesssim C_T U_t^3 L_t.
\]
Thus the simple linear Gronwall/localization argument from the two-layer note does not carry over automatically.
\end{remark}

\section{Implications}

The three-layer proxy is the first setting in which a shared $m\times m$ matrix appears, so it is the first reduced model that genuinely reflects the core parameter-efficiency mechanism of BatchEnsemble. At the same time, the concrete theorem becomes harder for a very specific reason:
\begin{enumerate}[leftmargin=2em]
  \item exact empirical closure still holds;
  \item the abstract bounded-Lipschitz large-$K$ theorem still applies if its hypotheses can be verified;
  \item what fails is the clean finite-horizon compact-region localization, because the natural majorant system is already quadratic at the middle layer and higher-order at the first layer.
\end{enumerate}
So the correct conclusion is not that the large-$K$ theory breaks in depth. Rather, the semidirect head mean-field structure remains plausible, but the two-layer proof strategy is no longer strong enough. A three-layer rigorous theorem would likely require one of the following:
\begin{enumerate}[leftmargin=2em]
  \item a sharper energy functional than the naive sum of layerwise norms;
  \item a one-sided Lipschitz or dissipativity structure for the deep shared dynamics;
  \item clipping or truncation at the concrete level;
  \item a short-time theorem plus continuation criterion.
\end{enumerate}

\paragraph{A heuristic hierarchy of head-specific parameters.}
The three-layer proxy also suggests a useful practitioner-facing distinction
between different locations of head-specific modulation.

First, the top-layer variables $a$ and $s_2$ act \emph{after} the last
nonlinearity. Conditional on the lower-layer features, the predictor is affine
in these variables, so centered coordinatewise perturbations are expected to
average out with width in much the same way as the readout coefficients do in
the two-layer note.

Second, the intermediate variables $s_1$ and $r_2$ act before averaging but are
themselves $m$-dimensional. If they are initialized with iid coordinates from a
common law and do not carry any additional low-dimensional structure, then each
head sees an empirical coordinate law that concentrates around the same
population law as $m\to\infty$. In that regime, these modulations may still
change the \emph{common} infinite-width predictor, but they are less natural
candidates for order-one head-to-head diversity that survives width.

By contrast, the input-side modulation $r_1\in\R^d$ is low-dimensional and
global: it affects every neuron coherently and does not self-average with $m$.
Thus if the infinite-width predictor map
\[
  r_1 \mapsto \Phi_t(r_1)\in L^2(\mathsf P_X)
\]
is nonconstant on the support of the head law, then $r_1$ is the most natural
source of width-persistent functional diversity. Likewise, nonlinearity of
$\Phi_t$ in $r_1$ is the most natural source of width-persistent ensemble bias.

This heuristic does \emph{not} say that middle-layer modulations can never
matter. Rather, it suggests a hierarchy:
\begin{enumerate}[leftmargin=2em]
  \item output-side and post-feature perturbations are the most likely to wash
  out with width;
  \item width-dimensional coordinatewise masks may alter the common limit but
  often self-average as sources of head-to-head variation;
  \item low-dimensional early modulations are the most plausible source of
  diversity and bias that remain visible even at very large width.
\end{enumerate}

For practitioners, the design implication is clear: if the goal is diversity
that survives increasing width, then it is more promising to place head-specific
parameters in early, low-dimensional, pre-nonlinear locations than to rely only
on late readout perturbations or width-dimensional coordinatewise masks.

\paragraph{Bottom line.}
Three layers are indeed more faithful and more practitioner-relevant than two, precisely because the shared hidden-to-hidden matrix makes efficient ensembling nontrivial. But this same feature is also the first place where the concrete proof stops being a straightforward extension of the two-layer argument.

\begin{thebibliography}{99}

\bibitem[Chen et~al.(2020)Chen, D'Amour, Kuleshov, and Shlens]{batchensemble}
Y.~Chen, A.~D'Amour, V.~Kuleshov, and J.~Shlens.
\newblock BatchEnsemble: An alternative approach to efficient ensemble and lifelong learning.
\newblock In \emph{International Conference on Learning Representations}, 2020.

\bibitem[Gorishniy et~al.(2024)Gorishniy, Kotelnikov, and Babenko]{tabm}
Y.~Gorishniy, A.~Kotelnikov, and A.~Babenko.
\newblock TabM: Advancing tabular deep learning with parameter-efficient ensembling.
\newblock arXiv:2410.24210, 2024.

\end{thebibliography}

\end{document}
